% RecSys 2026 — Short Paper (4 pages main + 2 pages refs)
% Double-blind, double-column ACM sigconf format
\documentclass[sigconf,anonymous,review]{acmart}

\usepackage{booktabs}
\usepackage{multirow}
\usepackage{tabularx}
\usepackage{mathtools}
\usepackage[capitalize,noabbrev]{cleveref}
\usepackage{xcolor}

\newcolumntype{Y}{>{\centering\arraybackslash}X}
\newcolumntype{L}{>{\raggedright\arraybackslash}X}


\setcopyright{none}
\copyrightyear{2026}
\acmConference[RecSys '26]{20th ACM Conference on Recommender Systems}{2026}{TBD}

\begin{document}

\title{Network Position as Confounder: Positional Debiasing and\\Variance-Weighted Learning for Graph-Based Uplift Modeling}

\begin{CCSXML}
<ccs2012>
<concept>
<concept_id>10002951.10003317.10003347.10003350</concept_id>
<concept_desc>Information systems~Recommender systems</concept_desc>
<concept_significance>500</concept_significance>
</concept>
<concept>
<concept_id>10010147.10010257.10010293.10010294</concept_id>
<concept_desc>Computing methodologies~Neural networks</concept_desc>
<concept_significance>300</concept_significance>
</concept>
<concept>
<concept_id>10002950.10003648.10003671</concept_id>
<concept_desc>Mathematics of computing~Probabilistic inference</concept_desc>
<concept_significance>300</concept_significance>
</concept>
</ccs2012>
\end{CCSXML}
\ccsdesc[500]{Information systems~Recommender systems}
\ccsdesc[300]{Computing methodologies~Neural networks}
\ccsdesc[300]{Mathematics of computing~Probabilistic inference}
\keywords{uplift modeling, individual treatment effect, graph neural networks, positional encoding, heteroscedasticity, causal inference}

\begin{abstract}
In graph-based uplift modeling, a node's structural position is an unobserved confounder that simultaneously drives treatment assignment bias and outcome noise heteroscedasticity---two error sources from the same root. We formalize this \emph{dual-root structure} through a PEHE decomposition and propose two plug-in corrections: \textbf{Graph Positional Encoding (GPE)} injects node2vec embeddings as a proxy confounder adjustment (bias reduction), while \textbf{Variance-Weighted Learning} downweights noisy samples (variance reduction). Neither module modifies the base model's architecture, making both applicable to any existing graph ITE estimator. Across four base models spanning 2016--2025, GPE reduces PEHE in 87\% of separate-head settings; combining both modules reduces the current SOTA's error by up to 29\% at high noise, confirming their complementarity. Code: \url{https://anonymous.4open.science/r/CAVIN}.
\end{abstract}

\maketitle

%% ====================================================================
\section{Introduction}\label{sec:intro}
%% ====================================================================

Estimating Individual Treatment Effects (ITE) on social networks is central to targeted recommendations: users in tightly-knit communities respond up to 2.3$\times$ more strongly to identical promotions than isolated users~\cite{luo2021behavioural}, a gap that i.i.d.-based estimators~\cite{kunzel2019metalearners, shalit2017estimating} cannot recover. A growing line of GNN-based methods addresses this---from NetDeconf~\cite{guo2020learning} (WSDM~2020) through GIAL~\cite{chu2021graph} and GNUM~\cite{zhu2023gnum} to the current SOTA GDC~\cite{hu2025gdc} (WSDM~2025)---yet each introduces a monolithic architecture whose innovations do not transfer to other models.

We argue that these methods all grapple with the same unrecognized root cause. A node's latent structural position $Z^*$---its community membership, centrality, and bridging role---simultaneously (1)~confounds treatment assignment ($Z^* \to T$ and $Z^* \to Y$), inflating estimation \emph{bias}, and (2)~drives heteroscedastic outcome noise ($Z^* \to \sigma^2$), inflating estimation \emph{variance}. This \emph{dual-root structure} means PEHE decomposes into a bias term and a variance term, each requiring a distinct correction. No prior graph ITE method targets both explicitly.

We propose two lightweight, complementary plug-in modules. \textbf{Graph Positional Encoding (GPE)} fuses node2vec embeddings with node features via cross-attention, implementing Veitch et al.'s proxy confounder theory~\cite{veitch2019using} to reduce the bias term. \textbf{Variance-Weighted Learning} estimates per-sample noise variance and downweights noisy pseudo-labels, reducing the variance term via Gauss--Markov efficiency~\cite{Aitken_1936}. GPE operates at the \emph{input} level; variance weighting at the \emph{loss} level. Neither modifies the base model's internal architecture.

Across four representative base models and five graph datasets, GPE reduces PEHE in 87\% of separate-head settings. At high noise, the combination reduces GDC's error by up to 29\%---confirming complementarity predicted by the decomposition. The R-learner~\cite{nie2021rlearner} provides theoretical precedent for variance-weighted pseudo-outcome regression; we extend this idea to the graph setting with learned heteroscedastic weights.


%% ====================================================================
\section{Problem Setup}\label{sec:prelim}
%% ====================================================================

Let $\mathcal{G} = (\mathcal{V}, \mathcal{E})$ be a graph where each node $i$ has features $X_i$, treatment $T_i \in \{0,1\}$, and outcome $Y_i$. Under standard network consistency, unconfoundedness, and overlap assumptions~\cite{guo2020learning}, the goal is to estimate $\tau(X_i) = \mathbb{E}[Y_i(1) - Y_i(0) \mid X_i, X_{\mathcal{N}}^{i}]$. Following the X-learner framework~\cite{kunzel2019metalearners}, we train outcome models $\mu_t$ per treatment arm and regress ITE estimates against cross-residuals $D_1 = Y(1) - \hat{\mu}_0$, $D_0 = \hat{\mu}_1 - Y(0)$.

\textbf{Dual-root structure.} Let $Z^*$ denote the latent structural position of a node (community membership, centrality, bridging role). We identify three simultaneous causal paths: $Z^* \to T$ (confounding), $Z^* \to Y(t)$ (heterogeneity), and $Z^* \to \sigma^2(Y)$ (heteroscedasticity). The PEHE decomposes accordingly:
\begin{equation}\label{eq:pehe_decomp}
    \mathrm{PEHE}^2 = \underbrace{\mathbb{E}\!\left[(\hat{\tau}(X) - \bar{\tau}(X))^2\right]}_{\text{bias}^2\;\text{(from confounding)}} + \underbrace{\mathbb{E}\!\left[\mathrm{Var}[\tilde{D} \mid X]\right]}_{\text{variance}\;\text{(from noise)}},
\end{equation}
where $\bar{\tau}(X) = \mathbb{E}[\tau \mid X]$. The bias term is inflated when $Z^*$ confounds $T \leftarrow Z^* \to Y$; the variance term when outcome noise $\sigma^2(X, Z^*)$ is position-dependent. Conditioning on a proxy for $Z^*$ reduces bias~\cite{veitch2019using}; inverse-variance weighting reduces the variance component via Gauss--Markov efficiency~\cite{Aitken_1936}. Critically, at high noise GPE alone can be counterproductive---the proxy amplifies position-dependent signal while that same position drives noise---so only the combination addresses both error sources.

%% ====================================================================
\section{Method}\label{sec:method}
%% ====================================================================

We present two plug-in modules---one at the input level, one at the loss level---that compose with any GNN-based ITE estimator without modifying its internal architecture.

\subsection{Graph Positional Encoding (GPE)}\label{sec:gpe}

GPE fuses raw node features $h \in \mathbb{R}^{n \times d_h}$ with pre-computed node2vec~\cite{grover2016node2vec} embeddings $p \in \mathbb{R}^{n \times d_p}$ via multi-head cross-attention:
\begin{align}
    Q &= hW_Q,\; K = pW_K,\; V = pW_V, \label{eq:qkv} \\
    \mathrm{head}_k &= \mathrm{softmax}\!\left(Q_k K_k^\top / \sqrt{d_k}\right) V_k, \label{eq:head}\\
    z &= \mathrm{LN}\!\bigl([h;\; W_o\,\mathrm{Concat}(\mathrm{head}_1,\ldots,\mathrm{head}_H)] + f_{\mathrm{res}}([h;p])\bigr), \label{eq:fuse}
\end{align}
where $W_Q, W_K, W_V, W_o$ are learnable projections and $f_{\mathrm{res}}$ is a residual projection. For any base model $M$, the GPE-enhanced version is $M_{\text{+GPE}}(\cdot) = M(\mathrm{GPE}(X, p), \cdot)$. Node2vec embeddings are computed once offline, adding negligible overhead.

\subsection{Variance-Weighted Learning}\label{sec:var}

We model $D_i = \tau(X_i) + \varepsilon_i$ with $\mathrm{Var}[\varepsilon_i \mid X_i] = \sigma^2(X_i)$. A dedicated variance head predicts $\log \sigma^2$ directly (guaranteeing positivity), trained on targets $\log((D_i - \hat{\tau})^2 + \epsilon)$ with $\hat{\tau}$ gradient-detached. Sample weights are:
\begin{equation}
    w_i = \frac{1}{\max\{\hat{\sigma}^2(X_i), \delta\}}, \quad \delta > 0,
\end{equation}
and the ITE loss becomes $\sum_{i} w_i (D_i - \hat{\tau}(X_i))^2$. For binary outcomes, $\sigma^2_t = \mu_t(1-\mu_t)$ is known analytically, requiring no learned head. The full training procedure is available in our code repository.

%% ====================================================================
\section{Experiments}\label{sec:exp}
%% ====================================================================

\subsection{Setup}

\begin{table}[t]
\caption{Graph dataset statistics. All use semi-synthetic treatment outcomes.}
\label{tab:datasets}
\centering
\footnotesize
\begin{tabular}{@{}lrrl@{}}
\toprule
Dataset & Nodes & Edges & Features \\
\midrule
CoraFull & 19,793 & 126K & 8,710 \\
DBLP & 17,716 & 106K & 1,639 \\
PubMed & 19,717 & 89K & 500 \\
BlogCatalog & 5,196 & 344K & 2,160 \\
Flickr & 7,575 & 12K & 12,047 \\
\bottomrule
\end{tabular}
\end{table}

We evaluate on five graph datasets (Table~\ref{tab:datasets}): three citation networks (CoraFull, DBLP, PubMed~\cite{bojchevski2018deep}) with community-aware semi-synthetic outcomes extending Guo et al.'s protocol~\cite{guo2020learning} (Louvain communities, multi-hop spillover, per-node heteroscedastic noise at $\rho \in \{5, 10, 30\}$), plus the standard BlogCatalog and Flickr benchmarks from the WSDM DGP~\cite{guo2020learning}. We test four representative base models: BNN~\cite{johansson2016learning} (S-learner), TARNet~\cite{shalit2017estimating} (T-learner), X-learner~\cite{kunzel2019metalearners}, and GDC~\cite{hu2025gdc} (WSDM~2025 SOTA). All share a 5-layer GAT backbone (dims [256,128,128,128,128], 4 heads). GPE uses 128-dim node2vec embeddings with 4-head cross-attention. Training: Adam ($\text{lr}{=}10^{-3}$), 200 epochs, 5 seeds; metric: rooted PEHE ($\downarrow$). Full results across seven models and additional ablations are at \url{https://anonymous.4open.science/r/CAVIN}.

\subsection{GPE as a Modular Plug-in}

Table~\ref{tab:main} reports rooted PEHE across three noise levels. GPE consistently reduces PEHE for all separate-head architectures (TARNet, X-learner, GDC) but not for BNN---the sole S-learner, whose shared $\mu(X,T)$ head cannot condition on position differently for treatment and control groups. Across all non-BNN settings in the full seven-model evaluation, GPE improves PEHE in 87\% of cells (Wilcoxon signed-rank $p < 10^{-5}$). TARNet benefits most: lacking built-in structural mechanisms, it has the most headroom for positional debiasing. GDC still gains substantially (DBLP $\rho{=}5$: $3.20 \to 2.76$, $-14\%$), indicating GPE captures context orthogonal to GDC's disentanglement. On the independently-generated BlogCatalog/Flickr benchmarks, GPE-enhanced TARNet also improves, ruling out DGP circularity.

\begin{table*}[t]
\caption{Rooted PEHE ($\downarrow$) with and without GPE. Mean$^{\pm\text{std}}$ over 5 seeds. \textbf{Bold} = +GPE improves. $^\star$GDC is the current SOTA~\cite{hu2025gdc}.}
\label{tab:main}
\centering
\footnotesize
\setlength{\tabcolsep}{3.0pt}
\begin{tabular}{@{}l l ccc ccc ccc@{}}
\toprule
& & \multicolumn{3}{c}{$\rho=5$} & \multicolumn{3}{c}{$\rho=10$} & \multicolumn{3}{c}{$\rho=30$} \\
\cmidrule(lr){3-5} \cmidrule(lr){6-8} \cmidrule(lr){9-11}
Model & & Cora & DBLP & PubMed & Cora & DBLP & PubMed & Cora & DBLP & PubMed \\
\midrule
\multirow{2}{*}{TARNet} & vanilla & $2.11^{\pm.22}$ & $1.88^{\pm.19}$ & $2.21^{\pm.10}$ & $3.01^{\pm.26}$ & $2.47^{\pm.31}$ & $3.48^{\pm.25}$ & $7.19^{\pm2.46}$ & $6.13^{\pm.55}$ & $4.19^{\pm.17}$ \\
& +GPE & $\mathbf{1.87}^{\pm.11}$ & $\mathbf{1.72}^{\pm.12}$ & $\mathbf{1.93}^{\pm.12}$ & $\mathbf{2.52}^{\pm.23}$ & $\mathbf{2.23}^{\pm.08}$ & $\mathbf{3.08}^{\pm.31}$ & $\mathbf{4.68}^{\pm.34}$ & $\mathbf{5.45}^{\pm1.20}$ & $4.36^{\pm.89}$ \\
\midrule
\multirow{2}{*}{X-learner} & vanilla & $2.30^{\pm.18}$ & $1.94^{\pm.12}$ & $2.12^{\pm.13}$ & $3.21^{\pm.49}$ & $2.42^{\pm.23}$ & $3.11^{\pm.10}$ & $4.37^{\pm.53}$ & $4.61^{\pm.90}$ & $3.93^{\pm.41}$ \\
& +GPE & $\mathbf{2.22}^{\pm.07}$ & $\mathbf{1.83}^{\pm.09}$ & $\mathbf{1.86}^{\pm.10}$ & $\mathbf{2.55}^{\pm.16}$ & $\mathbf{2.21}^{\pm.13}$ & $\mathbf{3.06}^{\pm.19}$ & $4.41^{\pm.50}$ & $\mathbf{4.49}^{\pm.66}$ & $\mathbf{3.91}^{\pm.10}$ \\
\midrule
\multirow{2}{*}{GDC$^\star$} & vanilla & $3.57^{\pm.40}$ & $3.20^{\pm.23}$ & $2.82^{\pm.10}$ & $4.03^{\pm.31}$ & $3.23^{\pm.34}$ & $3.45^{\pm.56}$ & $4.87^{\pm.26}$ & $4.20^{\pm.39}$ & $4.32^{\pm.28}$ \\
& +GPE & $\mathbf{3.09}^{\pm.34}$ & $\mathbf{2.76}^{\pm.31}$ & $\mathbf{2.56}^{\pm.10}$ & $\mathbf{3.49}^{\pm.60}$ & $\mathbf{2.72}^{\pm.11}$ & $\mathbf{2.93}^{\pm.42}$ & $5.59^{\pm1.71}$ & $\mathbf{3.79}^{\pm.19}$ & $4.44^{\pm.70}$ \\
\midrule
\multirow{2}{*}{BNN} & vanilla & $5.46^{\pm.46}$ & $4.98^{\pm.06}$ & $5.55^{\pm.06}$ & $5.41^{\pm.45}$ & $5.02^{\pm.03}$ & $5.44^{\pm.15}$ & $5.39^{\pm.46}$ & $4.98^{\pm.06}$ & $5.32^{\pm.13}$ \\
& +GPE & $5.50^{\pm.46}$ & $5.03^{\pm.06}$ & $5.54^{\pm.05}$ & $5.45^{\pm.44}$ & $5.02^{\pm.04}$ & $5.46^{\pm.07}$ & $5.46^{\pm.47}$ & $5.03^{\pm.04}$ & $5.44^{\pm.11}$ \\
\bottomrule
\end{tabular}
\end{table*}

\subsection{Variance Weighting: Complementing GPE}\label{sec:var_exp}

Table~\ref{tab:var} examines the interaction between GPE and variance weighting on GDC at $\rho=30$, where heteroscedasticity is most severe.

\begin{table}[t]
\caption{GDC ablation at $\rho=30$ (PEHE $\downarrow$, mean$^{\pm\text{std}}$, 5 seeds).}
\label{tab:var}
\centering
\footnotesize
\begin{tabular}{@{}lcccc@{}}
\toprule
Dataset & GDC & +GPE & +Var & +GPE+Var \\
\midrule
CoraFull & $4.64^{\pm.31}$ & $5.14^{\pm.95}$ & $5.04^{\pm.35}$ & $\mathbf{3.99}^{\pm.58}$ \\
DBLP & $4.02^{\pm.04}$ & $3.80^{\pm.14}$ & $3.87^{\pm.85}$ & $\mathbf{3.22}^{\pm.16}$ \\
PubMed & $4.47^{\pm.43}$ & $4.19^{\pm.60}$ & $\mathbf{3.19}^{\pm.34}$ & $3.55^{\pm.88}$ \\
\bottomrule
\end{tabular}
\end{table}

The results confirm the dual-root prediction: at high noise, GPE alone \emph{hurts} GDC on CoraFull ($4.64 \to 5.14$) because the positional proxy amplifies noise-correlated signal. Combining both modules yields the best PEHE on CoraFull and DBLP (up to 20\% reduction), while on PubMed variance weighting alone achieves a 29\% reduction---suggesting that when heteroscedasticity dominates confounding, noise suppression is the primary driver. On real-world tabular benchmarks, variance weighting yields a 4$\times$ Qini improvement on continuous outcomes (Hillstrom spend); full tabular results are in the supplementary materials.

%% ====================================================================
\section{Conclusion}\label{sec:conclusion}
%% ====================================================================

We identified network position as a dual error root in graph-based uplift modeling---simultaneously driving confounding bias and heteroscedastic variance---and proposed two plug-in corrections: GPE (input-level proxy confounder adjustment) and variance-weighted learning (loss-level noise suppression). Both modules compose with any existing GNN-based ITE estimator without architectural modification. GPE reduces PEHE in 87\% of separate-head settings; at high noise, the combination reduces GDC's error by up to 29\%. For practitioners, GPE is most effective on graphs with community structure and T-learner architectures; variance weighting adds the most value on continuous outcomes with heterogeneous noise.

\textbf{Limitations.} GPE is validated primarily on semi-synthetic graph outcomes; deployment-scale A/B testing would strengthen the practical case. Node2vec embeddings assume a static topology, limiting applicability to dynamic networks without periodic recomputation. Extended ablation studies (PE method variants, graph transformer comparison, design choices) are available in the supplementary materials.

\bibliographystyle{ACM-Reference-Format}
\bibliography{references}

\end{document}
